\section{The Sane Man believes in God}

\begin{quotex}
L'uomo sano crede in Dio e nella libertà del suo spirito. Anche senza rendersene conto, presuppone l'uno e l'altra, in ogni suo atto e in ogni sua parola. A cominciare dall'idea medieval, che è poi l'idea classica or greca, di Dio come \emph{esse quo maius cogitari nequit} o essere perfettissimo o assoluto, fino al più elevato concetto Cristiano di Dio come assoluto spirit che s'incarna nell'uomo e lo riscatta dalle sue tendenze naturali, la divinità è presente in tutti i suoi apetti essenziali dentro alla intelligenze e al cuore d'ogni uomo, come incrollabile realità, che resiste ad ogni dubbio e negazione, e perciò mai non fallisce, e regge tutte le cose e insieme il pensiero e la volontà degli uomini, quasi presupposto di ogni essere che esista e di ogni evento che accada.\\
\flright{\textsc{Giovanni Gentile}, \emph{Introduction to Philosophy}}
\end{quotex}


The sane man believes in God and in the freedom of his spirit. Even without realizing it, he presupposes them both, in his every act and in his every word. Beginning with the Medieval idea, which is also the classical or Greek idea, of God as \emph{the being of which nothing greater can be thought}, or the most perfect and absolute being, up to the most elevated Christian concept of God as the absolute spirit that is incarnated as man and redeems him from his natural tendencies, divinity is present in all his essential aspects within the intelligence and the heart of every man, as unshakable reality, that resists every doubt and negation, and therefore never fails, and supports all things and concurrently the thought and the will of men, as though the prerequisite of every being that exists and every event that occurs.


\flrightit{Posted on 2013-03-28 by Cologero}

\begin{center}* * *\end{center}

\begin{footnotesize}
\begin{sffamily}
\texttt{Logres on 2013-03-28 at 23:51 said:}

That's funny: I just got through musing aloud to my wife as to why some very ``noble'' and ``good'' people resist a belief in God: I suppose I conclude that somewhere, deep down, they know that would require a change or present a challenge or expose superstructures they would prefer not exposed. It's a very, very sore spot for them, and it continually amazes me how so many simple folk accept it, and so many intellectuals write pages and pages of narrative to show why morality and Being do not require the Christian God (or any other, presumably) for their existence. This is, in a sense true, in the sense that men can live in contradiction.

\hfill

\texttt{Cavalcare on 2013-03-29 at 19:14 said:}

Why should I `believe' in God? why should I not *know* Him instead? why this dogmatism, in relation to the Christian God (what is the ``Christian God''?)? there is more than one path to the peak of the mountain.

\hfill

\texttt{Cologero on 2013-03-29 at 22:46 said:}

You may be on to something, Cavalcare, or you may be missing the point. Although I have alluded to a project of translating two chapters of Evola, I didn't put this post in that context. We will see why Evola objected to Gentile's program, but first it is necessary to describe it. I also pointed out that Gentile's philosophy, in Evola's perspective, was close to being a way of realization, yet fell short. I can't cover all aspects in a comment, but I will quote asome Gentile, at the risk of being misunderstood or giving a false impression; they may even be closer to your point of view in some regards. By ``philosophy'', Gentile means what Guenon and Evola mean by metaphysics. However, unlike those two, Gentile relies totally on Western sources, being in the line of Plato, Aristotle, the Scholastics, etc. He falls short of them in that regard, although I think his understanding of society and political action has much to offer. Here goes:

\begin{quotex}\footnotesize

Religion is the philosophy of the multitudes and philosophy is the religion of the spirit, or, if you prefer, that of its highest representatives.

Catholicism is truly the most perfect religion, as modern European philosophy i.e., his is the most perfect philosophy. Together they are the highest creations of the Aryan spirit.
\end{quotex}


Perhaps you don't want to hear this, but what has surpassed it in the past 1700 years? Now Gentile was concerned with national identity which requires a spiritual basis; this is opposed to all individualism. Your personal revelations may be meaningful to you, but they have no bearing or influence on the world at large. If you adhere to that sort of spiritual individualism, I don't understand why you would care what others, who are likewise individuals, believer. Moreover, the issue is not what ``should'' be the spiritual tradition, but rather what actually has been the spiritual tradition. There are reasons for it, and we have provided plenty of them.

I will give you this one warning, Cavalcare. Comments need to refer to the post itself. The phrase ``Christian God'' was not used in the post. However, Gentile did offer a definition of God, which you can dispute if you like. Gentile mentioned that the Medieval idea of God is the same as the classical Greek understanding. So he is referring to classical theism, which apparently you need to research before offering your thoughts. Please do not confuse this with any misconceptions that you may have encountered, since we have no desire to bother to refute them.

\hfill

\end{sffamily}
\end{footnotesize}

